% Verbatim copy of /Users/alingalatan/workspaces/grothe/handoff/research-degree5/ai-notes/chatgpt-thread1-remote-cubic-comparison.tex (file dated 2026-11-05). Only this line was added.
% work-k/remote-cubic-comparison.tex
\documentclass[11pt]{article}
\usepackage[T1]{fontenc}
\usepackage{lmodern,amsmath,amssymb,mathtools,microtype}
\usepackage[a4paper,top=22mm,bottom=21mm,left=25mm,right=25mm]{geometry}
\usepackage{fancyhdr}
\usepackage[colorlinks=true,linkcolor=black,urlcolor=black]{hyperref}
\hypersetup{pdftitle={Remote cubics and a visible-dependent Gaussian term}}
\pagestyle{fancy}\fancyhf{}\renewcommand{\headrulewidth}{0pt}
\fancyhead[L]{\small\itshape Remote cubics and a visible-dependent Gaussian term}
\fancyfoot[C]{\small\thepage}\setlength{\headheight}{14pt}
\setlength{\parindent}{0pt}\setlength{\parskip}{4pt}
\setlength{\abovedisplayskip}{7pt plus 2pt minus 2pt}
\setlength{\belowdisplayskip}{7pt plus 2pt minus 2pt}
\newcommand{\E}{\mathbb E}
\newcommand{\PP}{\mathbb P}
\newcommand{\R}{\mathbb R}
\newcommand{\C}{\mathbb C}
\newcommand{\dd}{\,d}
\newcommand{\ind}{\mathbf 1}
\newcommand{\head}[1]{\section*{#1}}
\begin{document}
\thispagestyle{plain}
{\LARGE\bfseries Remote cubics and a visible-dependent\\[2pt] Gaussian term}\par
\medskip
Let \(T=|Z|^2\sim\operatorname{Exp}(1)\), \(a(t)=\alpha+\beta t\),
\(\alpha=\ell-2\beta\), \(r(t)=\sqrt t\,a(t)\), and \(L=Za(T)\). Write
\[
 F(G)=\E|L+G|+\E|L-G|,\quad M(v)=2\E|L+\sqrt v W|,\quad
 S=\sup_{v\ge0}\{M(v)-v/\varepsilon\},
\]
where \(W\) is an independent standard complex Gaussian. The box is
\(\varepsilon\in[0.74,0.76]\), \(\ell\in[0.76,0.90]\), \(\beta\in[0.12,0.16]\).
In particular \(\varepsilon\le\ell\), \(\alpha/\beta\ge11/4\).
We use the stated pure-remote case of Conjecture K when passing from a remote
comparison to the optimized inequality in Section 2.

\head{1. The missing term in the tangent bound}
Consider \(G=H+G_0\) with \(H\) remote and \(G_0(0,z')=0\). Write \(h=\|H\|_2^2\)
and use the constants \(C,D,E\) from Theorem 1 of \emph{Gaussian domination for
two classes of charge-one cubics}. Orthogonality and the phase of \(Z\) give
\[
 \E[|G|^2\mid T=t]=h+t\{C+D(t-2)^2+Et\},\qquad v=h+C+2D+2E .
\]
Let \(f(t)=(a(t)^2+v)^{-1/2}\), \(J=\int_0^\infty\sqrt t\,e^{-t} f(t)\dd t\).
In the square-root tangent estimate for \(F(G)/2\), the new term is
\[
 \frac h2 I_H,\qquad I_H=\int_0^\infty t^{-1/2}(1-t)e^{-t} f(t)\dd t
 =\int_0^\infty \sqrt t\,e^{-t}\{f(t)-2f'(t)\}\dd t>0.       \tag{1}
\]
The identity follows by integration by parts. In particular the proposed deletion
of the hypothesis on \(H\) fails at this step. Even with \(G_0=0\), the conditional
second-moment upper bound is \(\E\sqrt{|L|^2+h}>M(h)/2\) for \(h>0\), by strict
conditional Jensen for the Gaussian comparator.

There is a quantitative extension without any independence between \(H\) and
the coefficients of \(G_0\). Since \(-f'/f\le\beta/\alpha\), (1) gives
\(I_H\le(19/11)J\). The linear and quadratic estimates in the preceding theorem
bound the \(E,D\) coefficients by \(-J/2,-J/10\), respectively. Consequently
\[
 \boxed{\quad \frac{19}{11}h\le \frac{D}{10}+\frac E2
      \quad\Longrightarrow\quad F(H+G_0)\le M(v).\quad}       \tag{2}
\]
Here \(D\) is the squared coefficient of \(Z(T-2)\), and \(E\) the squared norm
of the linear driver of \(Z^2\bar V\). Pages 2--3 give another extension that
allows arbitrarily large \(h\).

\newpage
\head{An independent Gaussian scale mixture}
\textbf{Claim.} If \(G_0(0,z')=0\) and \(H\) is independent of \((Z,G_0)\) and
has the law of a centered circular Gaussian scale mixture, then
\(F(H+G_0)\le M(\|H\|_2^2+\|G_0\|_2^2)\) at all amplitudes.
For example \(H=XQ\) qualifies when \(X\) is standard complex Gaussian,
\(Q\) is an independent centered charge-zero quadratic, and the pair is independent
of the variables used by \(Z,G_0\).

First fix a deterministic Gaussian variance \(s>0\), set \(w=\|G_0\|_2^2\), and write
\(m_s(y)=\E|y+\sqrt s W|\) for \(y\ge0\), \(g_s(x)=m_s(\sqrt x)\). The Laplace
representation of the square root and a Gaussian integral give
\[
 g_s'(x)=\frac1{2\sqrt\pi}\int_0^{1/s} u^{-1/2}(1-su)^{1/2} e^{-ux}\dd u.       \tag{3}
\]
In fact the Laplace transform in the square-root representation is
\((1+sq)^{-1}\exp[-qx/(1+sq)]\); differentiate and set \(u=q/(1+sq)\).
Thus \(g_s\) is concave. Apply conditional Jensen to \(g_s(|L\pm G_0|^2)\)
and take its tangent at \(x(t)=t(a(t)^2+w)\). The two corrections, apart from
their nonnegative coefficients, are
\[
 \int t(t-2)e^{-t}g_s'(x(t))\dd t,\qquad
 \int t q(t)e^{-t}g_s'(x(t))\dd t,\qquad q(t)=t^2-4t+2 .
\]
The first is nonpositive because \(g_s'\circ x\) is decreasing and
\(\int t(t-2)e^{-t}\dd t=0\). For the second define
\[
 X(u)=\int_0^\infty t q(t)e^{-t-u x(t)}\dd t.
\]
Here \(X(0)=0\) and \(X'(0)<0\). Indeed \(x(t)\) has positive coefficients in
degrees one, two and three, and
\(\int t^{j+1}q(t)e^{-t}\dd t=j(j+1)(j+1)!\) for \(j\ge1\).
The function \(X\) changes sign just once on \(u>0\), from negative to positive.
Indeed let \(x_1,x_2\) be the values of \(x\) at \(2-\sqrt2,2+\sqrt2\).
If \(Y(u)=e^{ux_1}X(u)\), then \(e^{u(x_2-x_1)}Y'(u)\) is strictly increasing,
since its derivative integrates
\(tq(t)(x_1-x(t))(x_2-x(t)) e^{-t+u(x_2-x(t))}\ge0\).
Also \(Y(u)\to+\infty\): a compact interval inside \((0,2-\sqrt2)\) contributes
exponentially, whereas the negative part is bounded.

At \(s=0\) the preceding theorem proves \(\int_0^\infty u^{-1/2}X(u)\dd u\le0\).
Multiplication by the nonnegative decreasing weight \((1-su)_+^{1/2}\) preserves
this inequality: compare the weight on both sides of the unique sign change
to its value there. All interchanges are justified by absolute integrability
using \(g'_0(x)=1/(2\sqrt x)\). We have proved
\[
 \tfrac12 F(G_0+\sqrt s W)\le \E m_s\bigl(|Z|\sqrt{a(T)^2+w}\bigr).       \tag{4}
\]

\newpage
\head{The radial comparison after Gaussian smoothing}
For any nonnegative nondecreasing profile \(a\) of finite first moment
against \(|Z|\,d\PP\), with \(R=|Z|\), we have
\[
 \E m_s\bigl(R\sqrt{a(T)^2+w}\bigr)\le \E m_{s+w}(Ra(T)).                 \tag{5}
\]
For \(s=0\) this is the radial comparison in the preceding theorem. The case
\(w=0\) is an equality. Suppose \(s,w>0\). For constant profiles \(a=b\),
both sides equal \(\E R\sqrt{s+w+b^2}\). After differentiating in \(b>0\),
the difference of the right and left derivatives is the expectation of \(R\) times
\[
 m_{s+w}'(bR)-\frac b{\sqrt{b^2+w}}m_s'(R\sqrt{b^2+w}).                  \tag{6}
\]
This expression has at most one sign change, from negative to positive.
To verify this, \(m_s'(y)=m_1'(y/\sqrt s)\), and (3) gives
\[
 m_1'(y)=\mathrm{const}\, y\int_0^1 z^{-1/2}(1-z)^{1/2} e^{-y^2z}\dd z .
\]
Its logarithmic derivative with respect to \(\log y\) equals \(1-2\E_{\nu_{y^2}}x\),
where \(\nu_p\) has density proportional to
\(x^{-1/2}(1-x/p)^{1/2}e^{-x}\ind_{(0,p)}(x)\).
These laws increase stochastically with \(p\): for \(p_2>p_1\) the density
ratio increases on \((0,p_1)\), and the support extends to \(p_2\).
Thus \(m_1'(k_1R)/m_1'(k_2R)\) increases with \(R\) for \(0<k_1<k_2\).
Apply this with \(k_1=b/\sqrt{s+w}\), \(k_2=\sqrt{b^2+w}/\sqrt s\).

Since the full weighted integral of (6) is zero, its integral restricted to the
upper tail \(\{a(T)>b\}\) is nonnegative. Integrating in \(b\) proves (5).
Together with (4) it proves \(F(G_0+\sqrt s W)\le M(w+s)\).
Condition on the independent mixing variance and apply concavity of \(M\).
This proves the claim on page 2.

\textbf{A useful property of the visible law.} \(L\) itself is a centered
circular Gaussian scale mixture. For \(y\ge0\), let \(\psi(y)\) be the inverse of
\(t\mapsto t(\alpha+\beta t)^2\). Then \(\PP(|L|^2>y)=e^{-\psi(y)}\).
Differentiating in \(y\) is the operation \(g(t)\partial_t\) with
\(g(t)=[(\alpha+\beta t)(\alpha+3\beta t)]^{-1}\). Each factor of \(g\) is completely
monotone. Products and negative derivatives preserve complete monotonicity.
Starting with \(e^{-t}\) therefore proves that \(e^{-\psi(y)}\) is completely
monotone in \(y\). Bernstein's theorem represents this survival function as a
mixture of exponentials. The uniform phase gives the assertion about \(L\).
The representing measure has mass one and no atom at rate zero, by the values
at zero and infinity. See Schilling, Song and Vondra\v cek, \emph{Bernstein Functions},
Theorem 1.4. Also \(M\) is strictly concave since it is a mixture of
\(\sqrt\pi\sqrt{s+v}\) in \(v\).

\newpage
\head{2. An independent remote cubic and \((T-1)U\)}
\textbf{Theorem.} Let \(U\) be centered circular Gaussian, \(b=\E|U|^2\),
and let \(H\) be any square-integrable complex random variable independent of
\((Z,U)\), with \(U\) independent of \(Z\). If \(v=b+h\le\ell^2\), \(h=\E|H|^2\), then
\[
       \boxed{\quad F(H+(T-1)U)\le F(H+\sqrt b\,W).\quad}                  \tag{7}
\]
In fact moving the Gaussian variance from its constant value \(b\) to
\(b(T-1)^2\) along the affine segment decreases \(F\).

Write \(w_d(t)=b(1+d\,t(t-2))\) for \(0\le d\le1\) and
\(K_s(x)=\E|x+\sqrt s W|^{-1}\). Let \(\theta\) be an independent uniform angle.
The derivative of the interpolated \(F\) is
\[
 \frac b2\int_0^\infty e^{-t}t(t-2)
       \E_{\theta,H}K_{w_d(t)}(r(t)e^{i\theta}+H)\dd t.                   \tag{8}
\]
This follows from \(\partial_s m_s=K_s/4\) and averaging the phase of \(Z\).
For \(b>0\), differentiation is justified on compact subintervals of \(d<1\);
continuity gives the endpoint. We use the previously proved bound
\(K_s(x)\le6/(5|x|)\).

The positive and negative weights in (8) have equal mass \(N=4e^{-2}\).
The mean of \(t\) under the normalized negative weight on \((0,2)\) is
\(m_-=(9-e^2)/2<81/100\). There \(r(t)\le\ell\sqrt t\) and \(w_d(t)\le b\).
Jensen, first for the reciprocal radius and then in \(t\), bounds the magnitude
of the negative integral from below by
\[
       \frac{N}{\ell\sqrt{m_-+v/\ell^2}}>\frac{37N}{50\ell}.              \tag{9}
\]
For the positive integral fix the shift \(H\), set \(y=|H|/\ell\), and reduce
to radius \(\sqrt t\) as follows. If \(y\le\sqrt2\), set \(p=y\); increasing
the radius from \(\sqrt t\) to \(r(t)/\ell\) increases its distance to a point
at radius \(p\) for \(t\ge2\). If \(y>\sqrt2\), choose \(z\ge2\) with
\(r(z)/\ell=y\), and put \(p=\sqrt z\le y\). On radii at least \(\sqrt2\) the
map \(x\mapsto xa(x^2)/\ell\) increases both radii and their difference.
The squared-distance formula therefore gives in both cases
\[
       |r(t)e^{i\theta}+\ell y|\ge\ell|\sqrt t\,e^{i\theta}+p|,\qquad p\le y .
\]
The estimate on page 5 bounds the positive integral by
\[
       \frac{6N}{5\ell}\left(\frac12+\frac{7h}{60\ell^2}\right)
       \le \frac{37N}{50\ell}.
\]
Together with (9) this proves (7), strictly when \(b>0\).

\newpage
\head{The shifted outer integral}
Let
\[
 P(p)=\frac1N\int_2^\infty t(t-2)e^{-t}
             \E_\theta|\sqrt t\,e^{i\theta}+p|^{-1}\dd t .
\]
We prove \(P(p)\le\frac12+\frac7{60}p^2\) for all \(p\ge0\). First
\(P(0)=\frac14\E\sqrt{2+X}\le1/2\), \(X\sim\Gamma(2,1)\).
Put \(y=p^2\). By the binomial series and phase averaging the angular integral is
\[
 \max(t,y)^{-1/2}\sum_{n\ge0} a_n
 \left(\frac{\min(t,y)}{\max(t,y)}\right)^n,\qquad
 a_n=\left[\binom{2n}{n}/4^n\right]^2 .
\]
The logarithmic singularity at \(t=y\) is integrable; nonnegative terms permit
Tonelli. The \(n=0\) contribution is at most \(P(0)\). Since
\(\min(t,y)/\max(t,y)\le y/t\), the \(n=1\) contribution is at most
\(y\,\E(2+X)^{-1/2}/16\le y/(16\sqrt3)\). For the last inequality integrate
\(x/\sqrt{2+x}\) under an exponential law and apply concavity.

For \(n\ge2\), divide the contribution by \(y\) and set \(s=\log(t/y)\).
It is bounded by
\[
 a_n \sup_{t\ge2}\frac{(t-2)\sqrt t\,e^{-t}}N
       \left(\frac1{n-1}+\frac1{n+3/2}\right).
\]
Indeed on \(t\ge y\) the remaining factor is \(e^{-(n-1)|s|}\), and on \(t<y\)
it is \(e^{-(n+3/2)|s|}\). The supremum is below \(1/6\): with \(x=t-2\),
\(\sqrt{2+x}\le\sqrt3 e^{(x-1)/6}\), so \(x\sqrt{2+x}e^{-x}\le
(6\sqrt3/5)e^{-7/6}<2/3\).
The coefficient recurrence gives \(a_n\le1/(3n+1)\) by induction. Bounding by
\(1/(3n)\), replacing \(1/(n+3/2)\) by \(1/(n+1)\), and retaining the improvement
\(1/6-a_2=5/192\) at \(n=2\), the sum of the coefficients times the two fractions is
at most \(1/2-5/192\), by telescoping. Finally
\[
 \frac1{16\sqrt3}+\frac16\left(\frac12-\frac5{192}\right)<\frac7{60}.
\]
This proves the stated bound with elementary inequalities.

\textbf{Consequence for R2.} Along any direction of unit \(L^2\) norm, \(F(sG)\)
is \(2\)-Lipschitz in \(s\). Thus the penalized functional decreases for
\(s\ge\varepsilon\), and \(\varepsilon\le\ell\). When \(H\) is a remote cubic,
the stipulated pure-remote theorem applies also to \(H+\sqrt b W\) by a limit.
Use \(W_n=(2n)^{-1/2}\sum_{j=1}^n X_j(|X_j|^2-2)\) on fresh remote coordinates.
The planar central limit theorem gives \(W_n\Rightarrow W\), and bounded second
moments give convergence of the mean norms. The variances are exactly \(h+b\).
Consequently \(\Lambda(H+(T-1)U)\le S\) at every amplitude for disjoint coordinates.
This step uses the pure-remote theorem as an input; the fourth-moment upper bound
alone does not supply it.

\newpage
\head{3. Shared coordinates: small amplitudes}
Let \(G=(T-1)U+H\) with \((U,H)\) independent of \(Z\), where \(U\) is linear
and \(H\) a remote cubic. In particular \(\E[\bar U H]=0\), including for shared
coordinates. With \(b=\|U\|_2^2\), \(h=\|H\|_2^2\), \(v=b+h\), averaging the
phase gives, for a fixed direction with \(b>0\),
\[
 F(sG)-M(s^2v)=\frac{s^2 b}{2}
       \int_0^\infty \frac{t(t-2)e^{-t}}{r(t)}\dd t+o(s^2)<0
 \quad\text{for all sufficiently small }s>0.                              \tag{10}
\]
The integral is strictly negative because \(1/r(t)\) is strictly decreasing,
and the signed weight changes sign at 2 and has mass zero. The second-order
expansion follows by dominated convergence from the phase average of the norm;
its increment is bounded by \(|sG|^2/(2r(t))\).

\textbf{Uniform small amplitudes.} For every fixed \(\delta>0\), (10) extends
to a uniform same-variance theorem at sufficiently small \(s>0\) on all normalized
directions \(\|G\|_2=1\) with \(\E|G|^4\ge2+\delta\), throughout the box.
Together with the stated near-Gaussian result (ii), this gives the penalized
inequality for all such shared pairs at sufficiently small common amplitude.
The argument here gives an existence radius, not a radius of \(0.1\).

Here are details of the uniformity. In plane coordinates \(x=L\), write \(t(x)\)
for the inverse radial parameter and \(p(x)\) for the density of \(L\):
\[
 |x|^2=t(\alpha+\beta t)^2,\qquad
 p(x)=e^{-t}/\bigl[\pi(\alpha+\beta t)(\alpha+3\beta t)\bigr].
\]
The even function \(J(u,y)=\E_Z|L+(T-1)u+y|\) has a Taylor expansion in its four
real arguments through degree four with remainder bounded globally by
\(C(|u|+|y|)^6\), uniformly on the box. To justify this, \(t(x)\) is smooth by
the inverse function theorem, \(\nabla t\) is bounded, and its higher derivatives
are polynomially bounded, uniformly on the box. For small \(u\), the map
\(x\mapsto x+(t(x)-1)u\) has a smooth global inverse by contraction. Its Jacobian
is uniformly invertible. Differentiating the inverse and its Jacobian finitely
many times gives polynomial bounds in \(|x|\), since \(t\) and its derivatives
have such bounds. The pushed density and its parameter and spatial derivatives
through order six are consequently bounded by polynomials times
\(\exp[-c|x|^{2/3}]\) at infinity, uniformly for small \(u\). Translate the
pushed density by \(y\) before integrating against \(|x|\). This yields uniform
local \(C^6\) bounds. Evenness removes odd terms; linear growth of \(J\)
extends the remainder bound outside a neighborhood of zero.

Put \(B(y)=\E|L+y|\). Integration by parts gives
\[
 \Delta^2 B(0)=2\pi\int_0^\infty \frac{p(r)-p(0)}{r^2}\dd r<0 ,
\]
uniformly separated from zero on the compact box. Indeed \(\Delta|\cdot|=1/|\cdot|\)
and \(p\) strictly decreases. The pure-remote quartic Taylor coefficient of
\(F(sG)-M(s^2)\) is \(\Delta^2 B(0)(\E|H|^4-2)/32\); terms containing \(U\) in
the quartic coefficient are \(O(\sqrt b)\). All constants are dimension-free:
fixed-order Gaussian hypercontractivity bounds the sixth moments of \(U,H\)
uniformly, and \(\|(T-1)U\|_4=O(\sqrt b)\).

\newpage
\head{Shared coordinates: all-amplitude regions}
To finish the uniform argument, \(|\E|H|^4-\E|G|^4|=O(\sqrt b)\) by Hölder.
For small \(b\), the quartic coefficient is uniformly negative when the fourth
excess is at least \(\delta\). For \(b\) bounded below, the strictly negative
quadratic coefficient in (10) is uniform by compactness. The \(O(s^6)\) remainder
then proves the assertion.

\textbf{An explicit region for arbitrary sharing.} If \(\|H\|_2\le\|U\|_2/400\),
then \(\Lambda(G)\le S\) at all amplitudes, and \(F(G)\le M(v)\) when
\(\sqrt v\le\ell\). To prove it write \(A=(T-1)U\).
For \(h=0\), the positive bound on page 4 improves to \(3N/(5\ell)\).
Thus (8)--(9), integrated over \(d\), give
\[
         M(b)-F(A)\ge (7e^{-2}/25)b/\ell>(7/200)b/\ell.
\]
Let \(j(r,y)=\E_\theta|re^{i\theta}+y|\) for \(y\ge0\).
Then \(0\le j(r,y)-r\le5y^2/(16r)\). Indeed at \(r=1\) and \(y\le1\) use the
binomial series with positive squared coefficients and \(j(1,1)=4/\pi\).
For \(y\ge1\), symmetry and homogeneity bound \(j(1,y)\) by \(y+(4/\pi-1)/y\);
the minimum of \(y-y^2+5y^3/16\) on \([1,\infty)\) is \(8/27>4/\pi-1\).
Convexity gives \(j_y'\le5y/(4r)\) by taking the secant to \(2y\). Consequently
\[
 F(A+H)-F(A)\le\frac54\left(2\sqrt{bh}\,\E\frac{|T-1|}{r(T)}
                                   +h\,\E\frac1{r(T)}\right).
\]
Since \(\ell/\alpha\le19/11\), the two expectations times \(\ell\) are
less than \(5,7/2\), respectively, by bounding \(|T-1|\le T+1\).
With \(\sqrt h\le\sqrt b/400\) this error is below the displayed deficit.
The amplitude reduction on page 5 applies.

\textbf{The two examples.} Write \(X=z_2\), \(Y=z_3\), and
\[
 H_1=X(|X|^2-2),\qquad H_2=X(|Y|^2-1).
\]
For \(U=uX,\ H=cH_2\) with \(\Re(\bar u c)=0\), conditioning on \(Z,Y\) gives
a circular Gaussian with variance \(|u|^2(T-1)^2+|c|^2(|Y|^2-1)^2\).
This is also the conditional law obtained with independent Gaussian drivers
for the two summands, so R2 applies at all amplitudes.
For arbitrary phases, \(|c|\le|u|/16\) suffices: centering over \(Y\) and
\(|x+y|\le |x|+\Re(\bar x y)/|x|+|y|^2/(2|x|)\) bound the increase in \(F\)
by \(|c|^2\E[|X|^2/|L+u(T-1)X|]\). The inequality
\(s e^{-s}\le(2/e)e^{-s/2}\) and \(K_w(r)\le6/(5r)\) bound this by
\(24|c|^2\E[1/r(T)]/(5e)\), below the pure-\(A\) deficit for the indicated ratio.

\newpage
\head{What happens near a pure remote direction}
For \(i=1,2\), \(U=uX,\ H=cH_i\), there is also a constant \(\eta_i>0\)
such that \(|u|\le\eta_i |c|\) suffices for K at all amplitudes throughout
the box. This argument supplies no numerical value of \(\eta_i\).

First, both examples obey strict same-variance comparison for \(u=0,\ c\ne0\).
For \(H_2\) this is strict Jensen for \(M\) in the random variance
\(|c|^2(|Y|^2-1)^2\).
For \(H_1\), condition on the Gaussian mixing variance \(s\) of \(L\), as on
page 3. Concavity of \(g_s\), with tangent at \(2|c|^2t\), compares the expectations
of \(m_s(|c|\sqrt t|t-2|)\) and \(m_s(\sqrt2|c|\sqrt t)\).
The correction is proportional to
\(\int t q(t)e^{-t} g_s'(2|c|^2t)\dd t\). This is strictly negative:
in (3) the inner integral \(X(u)\) equals
\[
        \frac{2k(k-2)}{(1+k)^4},\qquad k=2|c|^2u.
\]
At \(s=0\) its integral against \(u^{-1/2}\) is strictly negative because
\(\int \sqrt t\,q(t)e^{-t}\dd t=-\Gamma(3/2)/4\).
The decreasing cutoff in (3) preserves strict negativity, also if the cutoff
vanishes before the sign change. This proves the strict comparison.

Next the two families admit a uniform small-amplitude ball for matched comparison,
including directions of small fourth excess. For normalized directions whose
\(b\) stays bounded below, use the quadratic term on page 6. As \(b\to0\) the
fourth moments converge to those of normalized \(H_1,H_2\), respectively \(44,18\),
by Hölder. The small-amplitude result on pages 6--7 therefore applies.
On the remaining compact annulus of amplitudes up to \(\varepsilon\), strict
comparison on the remote axis and uniform continuity give a positive angular
neighborhood of that axis. This proves the assertion about \(\eta_i\), including
larger amplitudes by penalized localization. Intermediate ratios and general
shared cubics remain unresolved here.

\textbf{Why the independent replacement fails even for these examples.}
Let \(p_s(r)=(t(r)-1)p(r)\) be a signed density and
\(D(y)=\int p_s(x)|x+y|\dd x\). All derivatives of \(p_s\) decay rapidly,
and \(p_s(r)>p_s(0)\) for \(r>0\). As on page 6,
\[
       \Delta^2D(0)=2\pi\int_0^\infty [p_s(r)-p_s(0)]r^{-2}\dd r>0 .
\]
The radial Taylor expansion is
\(\nabla D(y)=\Delta D(0)y/2+\Delta^2D(0)|y|^2y/16+O(|y|^5)\).
The remainder bound holds globally because derivatives of order six are bounded
(convolve derivatives of \(p_s\) with the bounded gradient of the norm).
At \(H=cH_i\), \(c>0\) real, differentiation in the real coefficient \(u\) at zero gives
\[
 \left.\tfrac12\partial_u F(cH_i+u(T-1)X)\right|_{u=0}
 =\frac{\Delta^2D(0)}{16}c^3\E[|H_i|^2\Re(\bar XH_i)]+O(c^5).
\]
The linear term vanishes by \(\E[\bar XH_i]=0\).
The cubic moments in this display are \(32\) and \(4\):
they are \(\E[S^2(S-2)^3]\) for \(S=|X|^2\), and
\(\E|X|^4\,\E[(|Y|^2-1)^3]\), respectively.
For each sufficiently small fixed \(c>0\) the derivative is positive, whereas
\(F(cH_i+uW)\) has derivative zero at \(u=0\). The R2 comparison thus strictly
reverses for small \(u>0\), even inside the application box. This does not imply
a reversal against \(S\).

\newpage
\head{Why marginal moment bounds cannot settle sharing}
There is a sharper distinction between moment information and the Wick structure.
Fix a box profile with a positive Gaussian optimizer \(\sigma_*\).
Such profiles exist analytically: at \((\varepsilon,\ell,\beta)=(0.76,0.76,0.16)\),
Jensen under \(\Gamma(1/2,1)\) gives
\(M'(0)\ge\sqrt\pi/[2(\alpha+\beta/2)]>1/\varepsilon\).

Take \(X\) standard complex Gaussian and \(S_X=|X|^2\).
Choose a sign \(\chi(S_X)\in\{-1,1\}\) and set \(H=\sigma_*\chi(S_X)X,\ U=X\).
Multiplying each circular shell by a sign preserves the Gaussian law of \(H\).
Let
\[
       k(s)=\frac{D'(\sigma_*\sqrt s)}{\sigma_*\sqrt s},
\]
using the radial derivative of \(D\) on page 8. The Taylor expansion and
\(\Delta^2D(0)>0\) show that \(k\) is nonconstant. Under the nonatomic
probability measure \(S_X\,d\PP\), take \(\chi=+1\) on the upper half of the
values of \(k\), and \(\chi=-1\) on the lower half, splitting a level set if needed.
Then \(\E[S_X\chi]=0\), while \(\E[S_X\chi k(S_X)]>0\). Hence \(U\) and \(H\)
are orthogonal in \(L^2\), \(\E|H|^4=2\sigma_*^4\), and
\[
 \left.\tfrac12\partial_u F(H+u(T-1)U)\right|_{u=0}
       =\sigma_*\E[S_X\chi k(S_X)]>0.
\]
At \(u=0\), \(\Lambda(H)=S\); the added penalty is \(u^2/\varepsilon\).
Thus the optimized comparison fails in this enlarged class for all sufficiently
small \(u>0\). This \(H\) has a Gaussian marginal and is a nonconstant radial
sign transform of \(X\), outside the cubic Wick chaos. The example explains why
even exact knowledge of a Gaussian remote marginal and its covariance with \(U\)
does not replace joint information from the chaos.

\medskip
\textbf{Scope.} The general matched-variance assertion in R1 with arbitrary remote
\(H\) and the general optimized assertion in R3 at intermediate amplitudes
and coefficient ratios are open in this argument. R2 is proved throughout the box,
with the stipulated optimized remote theorem as its only additional conjecture-specific input.
\end{document}
